\documentclass{article}

\usepackage[a5paper]{geometry}
\usepackage[spanish]{babel}
\usepackage{amsmath}
\usepackage{enumerate}
\usepackage{hyperref}
%\usepackage{tasks}
%\usepackage[inline]{enumitem}

\begin{document}

\pdfinfo{
   /Author (Luis Carrillo Guti\'errez)
   /Title  (Creating a PDF document using PDFLaTeX)
   /CreationDate (D:20040502195600)
   /Subject ('algebra)
   /Keywords (Centro de Esudios Preuniversitarios;Factorizacion;\'algebra, ciclo 2015-i,gu\'ia de estudio 5)
}

\begin{enumerate}

\item{Factorice el polinomio : $P(x) = 2x^5 + 3x^4 - 8x - 12$ y determine la cantidad de factores primos.
\begin{align*}
P(x) & = \underline{\mathsf{2x^5 + 3x^4}} - 8x - 12 \\ % factor comun
& = x^4(2x + 3) \underline{\mathsf{- 8x - 12}} \\
& = \mathsf{x^4(\underline{2x + 3}) - 4(\underline{2x + 3})} \\
& = \mathsf{(\underline{x^4 - 4})}(2x + 3) \\
& = (x^2 - 2)(x^2 + 2)(2x + 3) 
% & = (x - \sqrt{2})(x + \sqrt{2})(x^2 + 2)(2x + 3) 
\end{align*}
\begin{enumerate}
\item{4}
\item{2}
\item{5}
\item{\underline{3}}
\item{6}
\end{enumerate}
}

\item{Determine la suma de los factores primos del polinomio : \\ $Q(x) = 4x^2(x -7) - 3x(x-7) - x + 7$
\begin{align*}
Q(x) & = 4x^2(x -7) - 3x(x-7)\underline{\mathsf{ - x + 7}} \\
& = 4x^2(\underline{\mathsf{x -7}}) - 3x(\underline{\mathsf{x-7}}) - (\underline{\mathsf{x-7}}) \\
& = (\underline{\mathsf{4x^2 - 3x - 1}})(x - 7) \\
& = (4x + 1)(x - 1)(x - 7) \\
\Sigma & = (4x + x + x) + 1 - 1 - 7 = 6x - 7 
\end{align*}
\begin{enumerate}
\item{6x + 7}
\item{5x + 1}
\item{\underline{$6x - 7$}}
\item{$7x - 3$}
\item{4x + 7}
\end{enumerate}
}

\item{Un factor primo del polimonio : \\ $R(a;b;c) = a^2 - 3ab + ac + 2b^2 - 2bc$, es : 
\begin{align*}
R(a;b;c) & = a^2 - 3ab + ac + 2b^2 - 2bc \\ % ordenar
& = a^2 - 3ab + 2b^2 \underline{\mathsf{+ ac - 2bc}} \\ % factor comun
& = \underline{\mathsf{a^2 - 3ab + 2b^2}} + c(a - 2b) \\
& = a^2 - 3ab \underline{\mathsf{- ab}} + 2b^2 \underline{\mathsf{+ ab}} + c(a - 2b) \\
& = a^2 - 4ab + 2b^2 \underline{\mathsf{+ 2b^2}} + ab \underline{\mathsf{- 2b^2}} + c(a - 2b) \\
& = \underline{\mathsf{a^2 - 4ab + 4b^2}} + \underline{\mathsf{ab - 2b^2}} + c(a - 2b) \\
& = (\underline{\mathsf{a - 2b}})^2 + b(\underline{\mathsf{a - 2b}}) + c(\underline{\mathsf{a - 2b}}) \\
& = (a - 2b)(a - b + c)   
\end{align*}
\begin{enumerate}
\item{$a - b - c$}
\item{\underline{$a - 2b$}}
\item{$a + b + c$}
\item{$a - 2c$}
\item{$a + 2b + 3c$}
\end{enumerate}
}

\item{Factorice el polinomio : $P(x) = x^4 - x^3 + 64x - 64$, y \\ determine :
\begin{enumerate}[i)]
\item{La cantidad de factores primos.}
\item{La cantidad de factores primos lineales.}
\item{La cantidad de factores primos cuadr\'aticos.}
\end{enumerate}
\begin{enumerate}
\item{3; 1 y 2}
\item{4; 2 y 2}
\item{5; 2 y 3}
\item{3; 3 y 0}
\item{\underline{3; 2 y 1}}
\end{enumerate}
\begin{align*}
P(x) & = \underline{\mathsf{x^4 - x^3}} + \underline{\mathsf{64x - 64}} \\
& = x^3(\underline{\mathsf{x - 1}}) + 64(\underline{\mathsf{x - 1}}) \\
& = (\underline{\mathsf{x^3 + 64}})(x - 1) \\
& = (x + 4)(x^2 - 4x + 4^2)(x - 1) \\
& = (x^2 - 4x + 16)(x + 4)(x - 1)
\end{align*}
}

\item{Sea el polinomio : $P(x;y) = x^{15}y^8 - x^9y^{14}$, dadas las proposiciones :
\begin{enumerate}[I.]
\item{Tiene seis factores primos.}
\item{$x^2 + xy + y^2$, es un factor primo.}
\item{La suma de sus factores primos es $2x^2 + y^2$}
\end{enumerate}
\begin{enumerate}
\item{\underline{I y II}}
\item{I y III}
\item{Solo II}
\item{Solo III}
\item{II y III}
\end{enumerate}
\begin{align*}
P(x;y) & = \underline{\mathsf{x^{15}y^8}} - x^9y^{14} \\
& = x^9y^8(\underline{\mathsf{x^6 - y^6}}) \\
& = x^9y^8(x^3 - y^3)(x^3 + y^3) \\
& = x^9y^8(x + y)(x^2 - xy + y^2)(x^3 + y^3) 
% INCOMPLETO
\end{align*}
}

\item{Al factorizar : $Q(x) = (x^2 + 7x -3)^2 - 2(x^2 + 7x) + 3$; la suma de los t\'erminos independientes de los factores primos es :
\begin{enumerate}
\item{7}
\item{\underline{-8}}
\item{8}
\item{4}
\item{-4}
\end{enumerate}
}

\item{Factorice el polinomio : $R(x) = (x^2 + 1)^2 - 11x^2 - 401$ y obtenga la suma de los factores primos lineales.
\begin{enumerate}
\item{$x + 2$}
\item{$2x + 3$}
\item{$5x$}
\item{\underline{$2x$}}
\item{$4x$}
\end{enumerate}
}

\item{?`Cu\'antos factores primos lineales se obtienen al factorizar : $P(a) = a^5 - 3a^4 - a + 3$ ?
\begin{enumerate}
\item{5}
\item{2}
\item{1}
\item{4}
\item{\underline{3}}
\end{enumerate}
}

\item{Factorice : $P(x) = x^8 - x^4 + 22x^2 - 121$, e indique la suma de coeficientes de un factor primo.
\begin{enumerate}
\item{9}
\item{10}
\item{\underline{-9}}
\item{-11}
\item{13}
\end{enumerate}
}

\item{Al factorizar el polinomio : $Q(x) = abx^2 + a^2x + b^2x + ab$, un factor primo es:
\begin{enumerate}
\item{\underline{$ax + b$}}
\item{$a + x$}
\item{$b + x$}
\item{$x + ab$}
\item{$x + 1$}
\end{enumerate}
}

\item{Al factorizar : $M(x) = (x^2 - 5x + 3)^2 + 4x^2 - 20x + 15$ resulta $M(x) = (x-a)(x-b)(x-c)(x-d)$. Calcule $R = \sqrt{abc + d}$; siendo $a > b > c > d$ 
\begin{enumerate}
\item{4}
\item{\underline{5}}
\item{7}
\item{6}
\item{3}
\end{enumerate}
}

\item{Factorice : $P(x;y;z) = xyz + xy + yz + xz + x + y + z + 1$, y determine :
\begin{enumerate}[i)]
\item{El n\'umero de factores primos.}
\item{La suma de los t\'erminos independientes de los factores primos.}
\end{enumerate}
\begin{enumerate}
\item{\underline{3 y 3}}
\item{3 y 4}
\item{3 y 2}
\item{3 y 5}
\item{4 y 3}
\end{enumerate}
}

\item{Al factorizar : $P(x) = x^4 + 2x^3 - 10x^2 - 11x - 12$, la suma de factores primos lineales es :
\begin{enumerate}
\item{$4x$}
\item{$4x + 1$}
\item{\underline{$2x + 1$}}
\item{$2x + 3$}
\item{$2x$}
\end{enumerate}
}

\item{Factorice : $P(x) = x^5 + x^4 - 2x^3 - 2x^2 + x + 1$, e indique el n\'umero de factores primos.
\begin{enumerate}
\item{1}
\item{\underline{2}}
\item{3}
\item{4}
\item{5}
\end{enumerate}
}

\item{Factorice: $P(x) = x^4 -x^3 - 3x^2 - 3x -18$, e indique la suma de coeficientes de un factor primo.
\begin{enumerate}
\item{1}
\item{\underline{4}}
\item{5}
\item{6}
\item{0}
\end{enumerate}
}

\item{Factorice : $P(x) = x^5 + 4x^4 + x^3 - 6x^2$, e indique la suma de factores primos.
\begin{enumerate}
\item{\underline{$4(x + 1)$}}
\item{$4(x - 1)$}
\item{$4(x + 2)$}
\item{$4(x - 2)$}
\item{$4x$}
\end{enumerate}
}

\end{enumerate}
\end{document}
